\rok{Март 2026.}{А}

Трећи поправни тест из Алгоритама и структура података

Први практични тест одржан 06.03.2026.

\zadatak{1}{Завршавање имплементације BFS алгоритма}
Основна метода BFS алгоритма није потпуна. Допунити имплементацију ове методе која треба да омогући комплетно извршавање овог алгоритма.

\textbf{Додатне напомене за решавање задатка:}
\begin{itemize}
	\item Написати алгоритам.
	\item У Main методи креирати произвољан граф. Обезбедити начин за тестирање и проверу нове функционалности, односно позив BFS методе.
	\item У template-у је закоментарисана метода:
	\begin{Verbatim}[fontsize=\small]
public void BFS(int startVertex)
	\end{Verbatim}
\end{itemize}

\zadatak{2}{Заједнички карактери са провером позиције (Strict Intersection)}
Написати методу која прима два stringa. Из првог stringa треба исписати само оне карактере који се појављују у другом stringu, али под условом да се налазе на истом индексу (позицији) у оба stringa бар једном. Сваки такав карактер исписати само једном.

\textbf{Додатни захтеви за решавање задатка:}
\begin{itemize}
	\item Користити Hash табелу из шаблона за чување броја појављивања сваког карактера.
	\item Имплементацију писати у методи:
	\begin{Verbatim}[fontsize=\small]
public static void printCommonAtSameIndex(string str, string pattern)
	\end{Verbatim}
	\item Игнорисати разлику између малих и великих слова.
\end{itemize}

\textbf{Пример:}
\begin{itemize}
	\item \textbf{Улаз 1:} \texttt{str = "bAnana", pattern = "canaba"}
	\item \textbf{Излаз 1:} \texttt{"an"} -- само карактери \texttt{"a"}, \texttt{"n"} и \texttt{"a"} се налазе на истој позицији.
	\item \textbf{Улаз 2:} \texttt{str1 = "Kod", str2 = "Algoritam"}
	\item \textbf{Излаз 2:} \texttt{"Nema takvih karaktera"}
\end{itemize}

\zadatak{3}{Сума елемената у опсегу [k1, k2]}
Написати методу која израчунава и враћа суму свих вредности из BST стабла које се налазе у затвореном опсегу између \texttt{k1} и \texttt{k2}.

\textbf{Додатни захтеви за решавање задатка:}
\begin{itemize}
	\item Јавна метода: \texttt{public int sumRange(int k1, int k2)}.
	\item Приватна метода: \texttt{private int sumRange(Node current, int k1, int k2)}.
	\item Подразумева се да стабло има барем два чвора.
	\item Не улазити у подстабла која су сигурно ван опсега.
	\item Алгоритам \textbf{ОБАВЕЗНО} писати у оквиру класе \texttt{BinarySearchTree}.
\end{itemize}

\textbf{Потписи метода:}
\begin{Verbatim}[fontsize=\small]
public int sumRange(int k1, int k2)
private int sumRange(Node current, int k1, int k2)
\end{Verbatim}

\textbf{Примери:}

\begin{center}
\begin{minipage}{\textwidth}
\centering
\begin{forest}
for tree={
	circle,
	draw,
	thick,
	fill=gray!20,
	minimum size=8mm,
	inner sep=1pt,
	s sep=8mm,
	l sep=12mm,
	edge={-, thick},
	font=\small
}
[8
	[4
		[2
			[, phantom]
			[3]
		]
		[5]
	]
	[9
		[, phantom]
		[11]
	]
]
\end{forest}

\vspace{0.3cm}
\textbf{Слика 1.} Пример BST-а.
\end{minipage}
\end{center}

\begin{itemize}
	\item \textbf{Улаз 1:} \texttt{k1 = 5, k2 = 10}
	\item \textbf{Излаз 1:} 17 (8 + 9)
\end{itemize}

\begin{center}
\begin{minipage}{\textwidth}
\centering
\begin{forest}
for tree={
	circle,
	draw,
	thick,
	fill=gray!20,
	minimum size=8mm,
	inner sep=1pt,
	s sep=8mm,
	l sep=12mm,
	edge={-, thick},
	font=\small
}
[10
	[4
		[2
			[3]
			[, phantom]
		]
		[6
			[5]
			[9]
		]
	]
	[20
		[12
			[, phantom]
			[14]
		]
		[29
			[27]
			[, phantom]
		]
	]
]
\end{forest}

\vspace{0.3cm}
\textbf{Слика 2.} Пример BST-а.
\end{minipage}
\end{center}

\begin{itemize}
	\item \textbf{Улаз 2:} \texttt{k1 = 5, k2 = 14}
	\item \textbf{Излаз 2:} 37 (6 + 9 + 10 + 12)
\end{itemize}

\sablon{Шаблон трећег поправног теста март 2026. група А}{2025/Mart/GrupaA/AISP2025_PopravniTest3_Test2_GrupaA.linq}
